
\documentclass[10pt]{amsart}

\usepackage{amsmath,amssymb,amsthm}

\newtheorem{theorem}{Theorem}

\title{V033: Two-term asymptotics of Lanczos coefficients}
\author{Roberto Alfano}
\date{20 October 2026}

\begin{document}
\maketitle

\section*{Definitions}

Let $U_1,U_2,\ldots$ be independent random variables,
each uniformly distributed on $[-1/2,1/2]$, and put
\[
S_M=\sum_{\ell=1}^{M}U_\ell,
\qquad M\in\mathbb{N},\quad M\geq1.
\]

Let $\{P_j^{[M]}\}_{j\geq0}$ denote the monic orthogonal
polynomials for the probability distribution of $S_M$.
By symmetry, their three-term recurrence has the form
\[
P_{j+1}^{[M]}(x)
 =xP_j^{[M]}(x)
 -\beta_j^{[M]}P_{j-1}^{[M]}(x),
\]
where $P_0^{[M]}=1$, $P_1^{[M]}=x$, and $j\geq1$.

\begin{theorem}[Two-term Lanczos asymptotics]
For every fixed integer $j\geq1$,
\[
\boxed{
\beta_j^{[M]}
 =\frac{jM}{12}
 -\frac{j(j-1)}{20}
 +O_j(M^{-1})
}
\qquad (M\longrightarrow\infty).
\]
The implicit constant may depend on $j$.
\end{theorem}

\begin{proof}

Define the variance-one random variable
\[
Y_M=\sqrt{\frac{12}{M}}\,S_M.
\]
Its moment-generating function is
\[
\mathbb E e^{tY_M}
 =
\left(
\frac{\sinh\bigl(t\sqrt{12/M}/2\bigr)}
     {t\sqrt{12/M}/2}
\right)^M.
\]

The Taylor expansion
\[
\log\frac{\sinh z}{z}
 =\frac{z^2}{6}
  -\frac{z^4}{180}
  +O(z^6)
\]
gives, coefficientwise through every fixed power of $t$,
\[
\log\mathbb E e^{tY_M}
 =
\frac{t^2}{2}
 -\frac{t^4}{20M}
 +O(M^{-2}).
\]
Consequently,
\[
\mathbb E e^{tY_M}
 =
e^{t^2/2}
\left(1-\frac{t^4}{20M}\right)
+O(M^{-2}),
\]
where the final remainder is understood coefficientwise
in every fixed finite Taylor truncation.

Let $Z$ be a standard Gaussian random variable, and
let $H_n$ denote the monic probabilists' Hermite polynomial.
In particular,
\[
H_4(x)=x^4-6x^2+3.
\]

The Gaussian identity
\[
\mathbb E[H_4(Z)e^{tZ}]
 =t^4 e^{t^2/2}
\]
shows that, for every fixed polynomial $f$,
\[
\mathbb E[f(Y_M)]
 =
\mathbb E[f(Z)]
 -\frac{1}{20M}\mathbb E[H_4(Z)f(Z)]
 +O_f(M^{-2}).
\tag{1}
\]

Let $Q_n^{[M]}$ be the monic orthogonal polynomial of
degree $n$ for the distribution of $Y_M$, and set
\[
h_n^{[M]}
 =\mathbb E\!\left[Q_n^{[M]}(Y_M)^2\right].
\]

For fixed $n$, the moment matrices converge to the
positive-definite Gaussian moment matrices. Their
entries have expansions in powers of $M^{-1}$.
The linear systems defining the coefficients of the
monic orthogonal polynomials therefore give
\[
Q_n^{[M]}(x)
 =H_n(x)+M^{-1}R_n(x)+O_n(M^{-2}),
\]
coefficientwise, where $\deg R_n\leq n-1$.

By Gaussian orthogonality,
\[
\mathbb E[H_n(Z)R_n(Z)]=0.
\]
Thus (1) yields
\[
h_n^{[M]}
 =
\mathbb E[H_n(Z)^2]
 -\frac{1}{20M}
  \mathbb E[H_4(Z)H_n(Z)^2]
 +O_n(M^{-2}).
\tag{2}
\]

The Hermite product formula is
\[
H_n(x)^2
 =
\sum_{k=0}^{n}
 k!\binom nk^2 H_{2n-2k}(x).
\]
Using
\[
\mathbb E[H_n(Z)^2]=n!,
\qquad
\mathbb E[H_4(Z)^2]=4!,
\]
we obtain
\[
\mathbb E[H_4(Z)H_n(Z)^2]
 =6n(n-1)n!,
\]
including the zero cases $n=0,1$.

Substitution into (2) gives
\[
h_n^{[M]}
 =
n!\left(
1-\frac{3n(n-1)}{10M}
+O_n(M^{-2})
\right).
\]

For monic orthogonal polynomials, the recurrence
coefficient is the ratio of consecutive squared norms.
Hence the coefficient for $Y_M$ is
\[
\widetilde\beta_j^{[M]}
 =\frac{h_j^{[M]}}{h_{j-1}^{[M]}}
 =j\left(
1-\frac{3(j-1)}{5M}
+O_j(M^{-2})
\right).
\]

Finally, the change of variable
$Y_M=\sqrt{12/M}\,S_M$ implies
\[
\beta_j^{[M]}
 =\frac{M}{12}\widetilde\beta_j^{[M]}.
\]
Therefore,
\[
\beta_j^{[M]}
 =
\frac{jM}{12}
-\frac{j(j-1)}{20}
+O_j(M^{-1}),
\]
as claimed.
\end{proof}

\section*{Instructions for independent verification}

The reviewer should independently check:

\begin{enumerate}
\item The coefficientwise moment expansion and its
      $O(M^{-2})$ remainder for each fixed degree.
\item The perturbation of the monic orthogonal polynomials.
\item The Hermite product identity and squared-norm expansion.
\item The scaling of the recurrence coefficients.
\item The exact rational formulas for
      $\beta_5^{[M]},\ldots,\beta_{16}^{[M]}$
      supplied separately.
\end{enumerate}

No claim of novelty is inferred solely from the
existence of these twelve explicit coefficients.

\end{document}
